% content/c05-teilchen/de.tex — Kapitel 5: Welcher Hilbertraum? Das Teilchen auf der Geraden
% (deutsche Leitfassung)
%
% Stand 2026-11-11: ein Anfang. Ein Kern-Abschnitt trägt die beiden Abbildungen zu den
% Vorlesungen 7 und 8 und ihre Demos (docs/plan/lec78/DEMOS.md). Prosa und Formeln aus dem
% Skript v0.3: latex/lectures/de/vl07.tex (Positionsbasis, Ortsdetektor, Dichte) und vl08.tex
% (freies Teilchen, Gaußsches Wellenpaket, Impulsraum). Der übrige Stoff der beiden Vorlesungen
% (Impuls- und Energieoperator, Momente, drei Dimensionen, Quantisierungsregeln) folgt.
% Notation wie in vl08: Fourier-Transformierte mit Tilde (der Hut markiert Operatoren),
% Breiten \pwidth, \pwidth_t, \pwidth_p, Zerfließzeit \tspread.

\chapter{Welcher Hilbertraum? Das Teilchen auf der Geraden}\label{c05}

% Stimme: vl07.tex:5, 16-35.
Die Postulate aus Kapitel~\ref{c04} sagen nicht, welcher Hilbertraum, welcher
Hamiltonoperator und welche Observablen zu einem gegebenen System gehören. Für das wichtigste
System, das Teilchen auf der Geraden, ist der Hilbertraum $\Hil = \Ltwo{\mathbb{R}}$, der Raum
der quadratintegrierbaren Funktionen $\Psi\colon \mathbb{R} \to \mathbb{C}$ mit dem
Skalarprodukt $\braket{\Phi}{\Psi} = \int \Phi^*(x)\Psi(x)\,dx$. Man schreibt den Wert der
Funktion an der Stelle $x$ als $\Psi(x) = \braket{x}{\Psi}$ und rechnet mit den Kets
$\ket{x}$, als wären sie eine Orthonormalbasis, die \emph{Positionsbasis}; an die Stelle des
Kronecker-Deltas tritt die Diracsche Deltafunktion, $\braket{x}{y} = \delta(x-y)$, und aus
Summen über die Basis werden Integrale. Die Funktion $\Psi$ ist die Liste der Koeffizienten
von $\ket{\Psi}$ in der Positionsbasis; sie heißt die \emph{Wellenfunktion}.

\section{Ortsdichte und freies Teilchen}\label{c05.s1}

% Stimme: vl07.tex:40-66.
Ein Ortsdetektor überwacht ein Intervall $[a,b]$ und antwortet mit „ja“, wenn er das Teilchen
dort findet. Sein Operator ist
\begin{equation}\label{c05.eq.ortsdetektor}
  \op{\EffOp}_{[a,b]} = \int_a^b \ketbra{x}{x}\,dx .
\end{equation}
Im reinen Zustand $\ket{\Psi}$ spricht er mit der Wahrscheinlichkeit
\begin{equation}\label{c05.eq.intervall}
  \Prob\bigl[x\in[a,b]\bigr] = \melem{\Psi}{\op{\EffOp}_{[a,b]}}{\Psi}
  = \int_a^b \abs{\Psi(x)}^2\,dx
\end{equation}
an. Weil $\abs{\Psi(x)}^2$ nicht negativ ist, ist $p(x) = \abs{\Psi(x)}^2$ eine
Wahrscheinlichkeitsdichte, die Dichte der Ortsmessung. Die Wellenfunktion selbst ordnet jeder
Position eine \emph{komplexe} Zahl zu; der Versuchung, sich schon $\Psi$ als Dichte
vorzustellen, soll man widerstehen.

\begin{tpcfigure}{c05.fig.dichte}
  \demo{c05.demo.detektor}
  \caption{Dieselbe Wellenfunktion zweimal: das Gaußsche Wellenpaket
    $\Psi(x) = (2\pi)^{-1/4} e^{-x^2/4} e^{3ix}$. (a) Real- und Imaginärteil wechseln ständig
    ihr Vorzeichen; gepunktet die Einhüllende $\pm\abs{\Psi}$. (b) Vom Faktor $e^{3ix}$ bleibt
    in $p = \abs{\Psi}^2$ nichts übrig. Die Fläche über $[a,b]$ ist die Ansprechwahrscheinlichkeit
    des Detektors $\op{\EffOp}_{[a,b]}$, Gleichung~\eqref{c05.eq.intervall}. Der schmale Streifen
    der Breite $\varepsilon$ bei $c$ ist dieselbe Konstruktion, nur schmaler: Schrumpft
    $\varepsilon$, so schrumpft die Fläche mit.}
\end{tpcfigure}

% Stimme: vl07.tex:69-80.
Warum ein Intervall und nicht ein Punkt? Die Wahrscheinlichkeit, das Teilchen \emph{genau} an
der Stelle $x$ zu finden, ist immer null. Das Objekt $\ketbra{x}{x}$ ist für sich genommen kein
Operator: $(\ketbra{x}{x})^2 = \delta(0)\ketbra{x}{x}$, und $\delta(0)$ ist keine Zahl. Für
ein Intervall wird eines der beiden Integrale vom Delta erledigt, und
$\op{\EffOp}_{[a,b]}^2 = \op{\EffOp}_{[a,b]}$ ist ein Projektor.

% Stimme: vl08.tex:50-91.
Wie bewegt sich ein freies Teilchen der Masse $m$? Seine Schrödingergleichung
$i\hbar\,\partial_t\Psi = -\frac{\hbar^2}{2m}\,\partial_x^2\Psi$ wird von jeder ebenen Welle
$e^{i(kx - \hbar k^2 t/2m)}$ gelöst, und weil sie linear ist, auch von jeder Superposition:
\begin{equation}\label{c05.eq.freieloesung}
  \Psi(x,t) = \frac{1}{\sqrt{2\pi}}\int_{-\infty}^{\infty} \widetilde{\Psi}(k)\,
  e^{i\left(kx-\frac{\hbar k^2}{2m}t\right)}\,dk .
\end{equation}
Die Gewichte $\widetilde{\Psi}(k)$ legt die Anfangsbedingung fest: Bei $t = 0$ steht hier die
inverse Fourier-Transformation, und $\widetilde{\Psi}$ ist die Fourier-Transformierte der
Anfangswellenfunktion,
\begin{equation}\label{c05.eq.phik}
  \widetilde{\Psi}(k) = \frac{1}{\sqrt{2\pi}}\int_{-\infty}^{\infty} \Psi(x,0)\,e^{-ikx}\,dx .
\end{equation}

\begin{merksatz}[Das freie Teilchen in drei Schritten]
  Berechnen Sie die Fourier-Transformierte $\widetilde{\Psi}(k)$ der Anfangswellenfunktion
  $\Psi(x,0)$. Multiplizieren Sie jede Komponente $\widetilde{\Psi}(k)$ mit ihrer Phase
  $e^{-i\hbar k^2t/2m}$. Transformieren Sie zurück.
\end{merksatz}

Für das Gaußsche Wellenpaket
$\Psi(x,0) = (\pwidth\sqrt{2\pi})^{-1/2}\, e^{-(x-a)^2/4\pwidth^2}$ sind Hin- und
Rücktransformation Gaußsche Integrale. Das Betragsquadrat bleibt eine Glockenkurve um~$a$,
nur ihre Breite wächst:
\begin{align}
  \abs{\Psi(x,t)}^2 &= \frac{1}{\pwidth_t\sqrt{2\pi}}\;
  e^{-\frac{(x-a)^2}{2\pwidth_t^2}}, \label{c05.eq.zerfliessen}\\
  \pwidth_t &= \pwidth\sqrt{1+\frac{t^2}{\tspread^2}},
  \qquad
  \tspread = \frac{2m\pwidth^2}{\hbar}. \notag
\end{align}
Das Wellenpaket \emph{zerfließt}: Nach der Zeit $\tspread$ ist es um den Faktor $\sqrt{2}$
breiter, und je schmaler es anfangs war, desto kürzer ist $\tspread$.

% Stimme: vl08.tex:93-127.
Dieselbe Information steckt in der \emph{Impulswellenfunktion} $\Phi(p,t)$, den
Koeffizienten von $\ket{\Psi}$ in der Basis der ebenen Wellen zum Impuls $p = \hbar k$; zur
Zeit null ist $\Phi(p,0) = \widetilde{\Psi}(p/\hbar)/\sqrt{\hbar}$. Für das freie Teilchen ändert sich im
Impulsraum nur eine Phase,
\begin{equation}\label{c05.eq.impulszeit}
  \Phi(p,t) = e^{-\frac{it}{\hbar}\frac{p^2}{2m}}\,\Phi(p,0),
  \qquad
  \abs{\Phi(p,t)}^2 = \abs{\Phi(p,0)}^2 .
\end{equation}
Für das Gaußsche Wellenpaket der Breite $\pwidth$ ist $\abs{\Phi(p,0)}^2$ eine Gaußverteilung
um $p = 0$ mit der Standardabweichung $\pwidth_p = \hbar/2\pwidth$; das Produkt der beiden
Breiten ist zur Zeit null $\pwidth\,\pwidth_p = \hbar/2$.

\begin{tpcfigure}{c05.fig.fourier}
  \demo{c05.demo.zerfliessen}
  \caption{Das Gaußsche Wellenpaket zur Zeit $t = 0$, dreimal, mit den Breiten $\pwidth/2$,
    $\pwidth$ und $2\pwidth$, links im Orts-, rechts im Impulsraum; gleiche Linie heißt
    dasselbe Paket. Die Balken über den Kurven markieren $\pm$ eine Standardabweichung:
    Halbiert man die Breite im Ortsraum, so verdoppelt sie sich im Impulsraum. Die Reihenfolge
    der Breiten kehrt sich um, das Produkt bleibt. Mit der Zeit zerfließen die Kurven links nach
    Gleichung~\eqref{c05.eq.zerfliessen}, die rechts bleiben nach
    Gleichung~\eqref{c05.eq.impulszeit} stehen.}
\end{tpcfigure}
